% Appendix: full proofs of Proposition prop:criterion and of Proposition
% prop:compression / Theorem thm:root-error. Sources:
% research/projects/bie-enclosure/theory/port-subspace.md (eqs. (9)-(14),
% review 12/12 hold) and rayleigh-truncation.md (Prop. 1, Thm. 2, review
% 11 hold, 1 citation-scope note, 0 fail).
\appendix

\section{Proof of \texorpdfstring{\cref{prop:criterion}}{the count criterion}}\label{app:criterion}
Write $A=\mathsf E_{\mathrm L}$, $B=\mathsf E_{\mathrm R}$, let $V_\zeta=\diag(I,\zeta I)$, and let
$\Pi$ swap the two $K$-row blocks. Multiplying out the blocks of
\eqref{eq:pencil} and of $G_{\mathrm L},G_{\mathrm R}$ gives
\begin{equation}\label{eq:app-swap}
  (G_{\mathrm R}-\zeta G_{\mathrm L})V_\zeta=V_\zeta\Pi(A-\zeta B).
\end{equation}
For $\zeta\ne0$ take determinants and cancel $\det V_\zeta$. As a polynomial
identity this shows that the two generalized eigenvalue problems have the
same multipliers with the
same algebraic multiplicities, including $0$ and $\infty$. On the unit circle
$V_\zeta$ and $\Pi$ are unitary, so $G_{\mathrm R}-\zeta G_{\mathrm L}$ and $A-\zeta B$
have the same singular values there. In particular $\sigma_{\min}(W)\ge g>0$ for $W=G_{\mathrm R}-G_{\mathrm L}$
($\zeta=1$), and $W$ is invertible.

Let $H=(G_{\mathrm R}+G_{\mathrm L})W^{-1}$. Expanding the products gives
\begin{equation}\label{eq:app-skew}
  H^*\mathsf J+\mathsf JH=2W^{-*}FW^{-1},\qquad
  \Delta=\mathsf J^{-1}W^{-*}FW^{-1},\qquad
  (H-\Delta)^*\mathsf J+\mathsf J(H-\Delta)=0 .
\end{equation}
For $0\le t\le1$ let $H_t=H-t\Delta$, $R_t=\frac12(H_t+I)W$ and
$L_t=\frac12(H_t-I)W$. At $t=0$ these are $G_{\mathrm R}$ and $G_{\mathrm L}$. For
$|\zeta|=1$,
\begin{equation}\label{eq:app-homotopy}
  \bigl\|(R_t-\zeta L_t)-(G_{\mathrm R}-\zeta G_{\mathrm L})\bigr\|
  \le\frac{|1-\zeta|}{2}\,\|\Delta W\|
  \le\frac{\|F\|}{j\,\sigma_{\min}(W)}<\underline g\le g,
\end{equation}
using $\|\mathsf J^{-1}\|=1/j$. Hence $R_t-\zeta L_t$ is invertible on the
unit circle for every $t$. The number of zeros of
$f_t(\zeta)=\det(R_t-\zeta L_t)$ inside the circle is
$\frac1{2\pi\ii}\oint f_t'/f_t\,\dd\zeta$. It depends continuously on $t$ and is
an integer, so it is constant. A multiplier may pass through $\infty$ along
the homotopy; this does not affect the count inside the circle.

At $t=1$, \eqref{eq:app-skew} shows that $H_1$ is similar to $-H_1^*$, so its
eigenvalues come in pairs $h\leftrightarrow-\bar h$ with equal algebraic
multiplicities. Moreover
\[
  R_1-\zeta L_1=\tfrac12\bigl[(1-\zeta)H_1+(1+\zeta)I\bigr]W .
\]
An eigenvalue of $H_1$ on the imaginary axis would give a multiplier on the
unit circle, which \eqref{eq:app-homotopy} excludes. So $H_1$ has $K$
eigenvalues in each open half-plane. The map $\zeta=(h+1)/(h-1)$ sends the
left half-plane into the unit disk and the right half-plane to its exterior,
with $h=1$ going to $\infty$; by the factorization above it preserves
multiplicities. With the homotopy and \eqref{eq:app-swap} the claim follows.
\qquad$\square$

\begin{remark}[a computable $\underline g$]\label{rem:app-gbar}
Let $\zeta_1,\dots,\zeta_N$ be equispaced on the unit circle and
$\ell_j\le\sigma_{\min}(A-\zeta_jB)$. Since $\sigma_{\min}$ is Lipschitz with
constant $\|B\|$ in $\zeta$ and every point of the circle lies within the chord
$2\sin(\pi/(2N))$ of a node,
$\underline g=\min_j\ell_j-2\sin(\pi/(2N))\,\|B\|$ is a lower bound for $g$.
The experiments of this paper report observed counts and flux defects and
do not evaluate this bound.
\end{remark}

\section{Proofs for the Rayleigh truncation}\label{app:root-error}
We use the setting of \cref{sec:truncation}. In addition to
\cref{ass:truncation} we use that every cell problem of the truncated model is
uniquely solvable (\cref{ass:cell}), and that the domain of $A(\beta)$ is
$H^2_\beta(B)$; both hold for the piecewise-constant coefficients and $C^2$
interfaces of \cref{sec:geometry}. \todo{Cite a published regularity statement
for $D(A)=H^2_\beta(B)$.}

\subsection{Proof of \texorpdfstring{\cref{prop:compression}}{the compression proposition}}
\emph{One deleted order.} In a gap strip $G_j$ put $t=x-l_j$ and $w=w_j$, and
let $\kappa=\kappa_m(\lambda)$, with the branch positive on the real axis and
$\Re\kappa>0$ on $U_r$. If $(\alpha,\delta)$ are the values of the $m$th Fourier
coefficient at $t=0$ and $t=w$, the full model gives
\[
  f(t)=\frac{\sinh(\kappa(w-t))}{\sinh(\kappa w)}\,\alpha
       +\frac{\sinh(\kappa t)}{\sinh(\kappa w)}\,\delta,\qquad
  \begin{bmatrix}-f'(0)\\ f'(w)\end{bmatrix}=D_m\begin{bmatrix}\alpha\\ \delta\end{bmatrix}.
\]
In the truncated model the rods on the left of the strip see only their own
outgoing wave $\alpha\ee^{-\kappa t}$ in this order, and the rods on the right
only $\delta\ee^{-\kappa(w-t)}$. The outward derivatives at the two ends are
then $\kappa\alpha$ and $\kappa\delta$, wherever the artificial boundary or wall lies in the strip.
Truncation therefore replaces $D_m$ by $\kappa I_2$, which gives
\eqref{eq:strip-defect}. If a strip contains several such lines, the empty cells
between them carry no incident field in a deleted order and radiate in both
directions, so this order vanishes there, and the rods at the two ends see the
same $\kappa$.

\emph{Trace bound.} Let $h_m=(1+\beta_m^2)^{1/2}$ and $s=\min(w,h_m^{-1})$. For
$f\in H^1(0,w)$, averaging $f(0)=f(x)-\int_0^xf'$ over $0<x<s$ gives
\[
  h_m|f(0)|^2\le2h_ms^{-1}\int_0^s|f|^2+2h_ms\int_0^s|f'|^2
  \le2\max(1,w_{\min}^{-1})\int_0^w\bigl(h_m^2|f|^2+|f'|^2\bigr),
\]
since $h_ms\le1$ and $h_m/s\le\max(1,w_{\min}^{-1})h_m^2$. The same holds at
$t=w$. Apply this to $f=u_m$, sum over $m$ (orthogonality of $\psi_m$), use
$\rho=\rho_{\mathrm e}$ in the strip, and sum over the disjoint strips:
\begin{equation}\label{eq:app-trace}
  \sum_j\sum_mh_m\|t_{jm}u\|^2\le C_{\mathrm{tr}}^2\|u\|_V^2 .
\end{equation}

\emph{Size of $\Delta_m$.} On $U_r$, $\Re\kappa_m\ge q_M(r)$ for $|m|>M$, because
$(\Re\sqrt z)^2=(|z|+\Re z)/2\ge\Re z$ for $\Re z>0$; and $|\kappa_m|\le C_\kappa h_m$.
The symmetric matrix $S$ has the orthonormal eigenvectors $(1,\pm1)/\sqrt2$, so
$\Delta_m$ is unitarily diagonalizable with eigenvalues $2\kappa p/(1+p)$ and
$-2\kappa p/(1-p)$, $p=\ee^{-\kappa w}$. With $|p|\le\ee^{-q_M(r)w_{\min}}$,
\begin{equation}\label{eq:app-delta}
  \sup_{j,\,|m|>M}\frac{\|\Delta_m(w_j,\lambda)\|_2}{h_m}
  \le\frac{2C_\kappa}{\ee^{q_M(r)w_{\min}}-1}.
\end{equation}
\eqref{eq:app-trace} and \eqref{eq:app-delta} applied to both traces in the
definition of $\mathcal E_M$ give \eqref{eq:compression-bound}. Each
$\Delta_m/h_m$ is holomorphic and uniformly bounded on a slightly larger disk.
By Cauchy's estimates the block-diagonal multiplier is holomorphic in
operator norm, and so is $\mathcal E_M$. For real $\lambda$ each $\Delta_m$ is real
symmetric, and $\Delta_m(w,\bar\lambda)=\Delta_m(w,\lambda)^*$.

\emph{Kernel and truncated modes.} Let $\lambda$ be real. Given a truncated mode,
keep the field in every block containing rods and the retained orders in the
strips, and fill the deleted orders of each strip by the full-model
interpolant above from the two end values. The filled field has finite energy,
since $\int_0^w(|f'|^2+\kappa^2|f|^2)=v^*D_mv\le\kappa\coth(\kappa w/2)\|v\|^2$ for
$v=(\alpha,\delta)$ and $h_m^2/\kappa_m^2$ is bounded for the deleted orders. Its
weak form differs from that of the truncated model exactly by
$\sum(t_{jm}v)^*(\kappa I_2-D_m)t_{jm}u$, so the field lies in
$\ker\mathcal T_M(\lambda)$. Conversely, for $u\in\ker\mathcal T_M(\lambda)$, test
functions vanishing at the strip ends show that each order solves its
ordinary differential equation in the strip. Replace the deleted orders by
the two one-sided outgoing exponentials with the same end values. The weak
form then gives the matching of normal derivatives at the strip ends, which is
the outgoing condition of the truncated model for the deleted orders. By
uniqueness of the cell problems this is a truncated mode. The two
constructions are inverse to each other. \qquad$\square$

\subsection{Proof of \texorpdfstring{\cref{thm:root-error}}{the root-error theorem}}
\emph{Complement.} Let $\phi$ be the eigenfunction of $A(\beta)$ at $\lambda_*$
with $\|\phi\|_V=1$, $P$ the $V$-orthogonal projection onto $\CC\phi$, and
$Q=I-P$. The map $u\mapsto(A+I)^{1/2}u$ is unitary from $V$ onto
$L^2(B,\rho)$ and turns $\mathcal T(\lambda)$ into $(A-\lambda)(A+I)^{-1}$. Hence
$P$ and $Q$ reduce $\mathcal T(\lambda)$, the $P$ block is the scalar
$s(\lambda)=(\lambda_*-\lambda)/(1+\lambda_*)$, and for $\lambda\in U_r$ and
$t\in\sigma(A)\setminus\{\lambda_*\}$,
\[
  \frac{1+t}{|t-\lambda|}\le1+\frac{|1+\lambda|}{|t-\lambda|}\le1+\frac{1+\lambda_*+r}{g-r}=G .
\]
So $\|(Q\mathcal T(\lambda)|_{QV})^{-1}\|\le G$, and by a Neumann series and
\eqref{eq:M-condition}
\[
  H_M(\lambda)=\bigl(Q\mathcal T_M(\lambda)|_{QV}\bigr)^{-1},\qquad
  \|H_M(\lambda)\|\le\frac{G}{1-G\varepsilon_M}\le2G .
\]

\emph{Schur complement.} Write $u=c\phi+v$ with $v\in QV$. The $Q$ rows give
$v=-cH_MQ\mathcal E_M\phi$, so $\mathcal T_M(\lambda)$ has a nontrivial kernel iff
\[
  s_M(\lambda)=s(\lambda)+(\mathcal E_M\phi,\phi)_V
  -(\mathcal E_MH_MQ\mathcal E_M\phi,\phi)_V=0 .
\]
$s_M$ is holomorphic on $U_r$ and
\[
  \sup_{U_r}|s_M-s|\le\varepsilon_M+\frac{G\varepsilon_M^2}{1-G\varepsilon_M}
  =\frac{\varepsilon_M}{1-G\varepsilon_M}\le2\varepsilon_M .
\]
On $\partial U_r$, $|s|=r/(1+\lambda_*)>8\varepsilon_M>2\varepsilon_M$ by
\eqref{eq:M-condition}. By Rouché's theorem $s_M$ has exactly one zero
$\lambda_M$ in $U_r$, counted with multiplicity. Since
$\mathcal E_M(\bar\lambda)=\mathcal E_M(\lambda)^*$, also
$s_M(\bar\lambda)=\overline{s_M(\lambda)}$, so the zero is real. Evaluating
$|s(\lambda_M)|=|s(\lambda_M)-s_M(\lambda_M)|\le2\varepsilon_M$ gives
$|\lambda_M-\lambda_*|\le2(1+\lambda_*)\varepsilon_M<r/4$. Then
$k_M>(\lambda_*-r)^{1/2}$ and
$|k_M-k_*|=|\lambda_M-\lambda_*|/(k_M+k_*)$ give \eqref{eq:root-error}.
\Cref{prop:compression} turns $\lambda_M$ into a truncated mode.

\emph{Rate at $k_*$.} By the definitions of $q_M$ and $q_M(r)$,
\[
  0\le q_M-q_M(r)=\frac{\rho_{\mathrm e}r}{q_M+q_M(r)}
  \le\frac{\rho_{\mathrm e}r}{q_{M_0}+q_{M_0}(r)},
\]
so the exponent $-q_M(r)w_{\min}$ in \eqref{eq:root-error} may be replaced by
$-q_Mw_{\min}$ at the cost of a constant independent of $M$. Nothing in the
argument depends on the positions of the artificial boundaries and walls inside their strips.
\qquad$\square$

\todo{Proofs of the half-strip paragraph, the first-order formula
\eqref{eq:first-order} and the $H^2$ refinement are in the theory note
(sections 3 and 4.2); decide whether to include them here or in a
supplement.}
